\documentclass[12pt,a4paper]{article}
\usepackage[utf8]{inputenc}
\usepackage[T5,T1]{fontenc}
\usepackage[a4paper,margin=2.5cm]{geometry}
\usepackage{setspace}
\usepackage{booktabs}
\setstretch{1.12}

\title{\textbf{Personal Progress Report}}
\author{{\fontencoding{T5}\selectfont Ngô Thảo Anh} -- 25125003\\
\small SC203 -- Group 8}
\date{}

\begin{document}
\maketitle

\begin{abstract}
\noindent SonicMaster is a model that restores and masters music based on text instructions, such as a request to reduce room echo. This report covers the reproduction on 7,000 test clips, a pilot on preserving audio detail, an analysis of joins in full tracks, and a listening review of 100 clips. The overall scores were close to those in the paper. The pilots revealed trade-offs between measured similarity and quality, while listening identified recurring issues such as buried sounds and robotic vocals.
\end{abstract}

\section{Weekly Objectives}
SonicMaster first compresses the audio with a variational autoencoder (VAE). The model edits this compressed representation based on the instruction, and the VAE decoder turns the result back into audio \cite{sonicmaster}. My goals were to reproduce the published results, test whether adding back audio detail helps, check whether chunk joins introduce loudness errors, and identify recurring audible issues.

\section{Work Completed}
\begin{itemize}
  \item Completed inference and evaluation for 7,000 official test clips from 1,000 source tracks. The dataset combined 1,400 existing outputs, 14 benchmark clips, and 5,586 production clips run on Kaggle. I compared the results with Tables 1--3 of the paper.
  \item Ran the VRG pilot and compared it with a simple blend of the input and the SonicMaster output.
  \item Checked joins in 12 full tracks using two sets of chunk boundaries and compared them with other regions in the same tracks.
  \item Reviewed 100 clips by listening and recorded the timing and type of audible issues.
\end{itemize}

\section{Experiments and Findings}

\subsection{Reproduction on 7,000 Test Clips}
I used the published weights, 10 Euler steps, guidance scale 1, seed 0, and FP32. The table shows the overall scores for SonicMaster \cite{reproduce}.

\begin{center}
\begin{tabular}{lrr}
\toprule
Metric & Reproduction & Paper \\
\midrule
FAD-CLAP $\downarrow$ & 0.07256 & 0.073 \\
KL ($\times 1{,}000$) $\downarrow$ & 0.89278 & 0.888 \\
Mel-SSIM $\uparrow$ & 0.60907 & 0.609 \\
PQ $\uparrow$ & 7.70506 & 7.705 \\
\bottomrule
\end{tabular}
\end{center}

Compared with the input, SonicMaster improved average KL and PQ, while Mel-SSIM fell from 0.777 to 0.609. This shows that the overall spectral balance and estimated production quality can improve while the detailed spectrogram structure becomes less similar to the clean reference. Average error fell for all 19 measured degradation attributes. The clipping score differed substantially from the paper and still needs a check of the scoring method.

\subsection{Adding Back Detail with VRG}
VRG adds a fraction of the difference between the input and its VAE reconstruction to the SonicMaster output. A weight of 0.1 raised Mel-SSIM by 0.0098, below the preset threshold of 0.01. At a weight of 0.25, the error for one EQ attribute exceeded the allowed limit.

A comparison blend of 90\% SonicMaster output and 10\% input raised Mel-SSIM by 0.0125 on 140 additional clips from 20 tracks. Average PQ fell by 0.0136 points, and the change in FAD varied by degradation group. The blend improved spectrogram similarity, with mixed results on other measures \cite{pilots}.

\subsection{Checking Joins in Full Tracks}
Each 30-second chunk was about 46 milliseconds shorter after VAE decoding. Joining chunks based on their decoded length can cause timing drift, so I checked a method that places each chunk at its original input timestamp. On the 12 tracks, the join regions had no larger loudness errors than their comparison regions under the chosen measure. The full-track PQ estimate was lower than the input for 11 of 12 tracks. The pilot did not establish the cause of that drop \cite{pilots}.

\subsection{Listening Review of 100 Clips}
I listened to five 6-second excerpts per clip. Of the 100 clips, 49 were labeled as containing vocals and 51 as instrumental. I recorded at least one issue in 68 clips. The most frequent categories were buried, missing, or less distinct sounds (39 clips) and robotic or Auto-Tune-like vocals (38 clips). Categories could overlap.

In a follow-up comparison of 10 vocal excerpts, I first heard the robotic effect in the VAE reconstruction in the cases checked. KL and PQ could also improve on excerpts where an instrument remained hard to hear \cite{listening}.

\section{Ongoing Work and Limitations}
\begin{itemize}
  \item The clipping discrepancy needs a comparison of the scoring formula and inputs on the same clips.
  \item The join pilot measured loudness on 12 selected tracks. The VRG pilot used small sets, and its scores do not establish a listening benefit.
  \item The listening results describe the 100 clips reviewed. Further annotation is needed to distinguish buried instruments from buried vocals and identify issues already present in the input.
\end{itemize}

\section{Plan for Next Week}
\begin{enumerate}
  \item Continue identifying and documenting audible issues in SonicMaster's outputs.
  \item Review the literature on music restoration and summarize work relevant to the observed issues.
\end{enumerate}

\begin{thebibliography}{9}
  \bibitem{sonicmaster} J. Melechovsky et al. \textit{SonicMaster: Towards Controllable All-in-One Music Restoration and Mastering}. ICML, 2026.
  \bibitem{reproduce} SC203 Group. \textit{SonicMaster Reproduction Report on 7,000 Test Clips}. 11 November 2026.
  \bibitem{pilots} SC203 Group. \textit{VRG Pilot Report} and \textit{Full-Track and Chunk Join Pilot}.
  \bibitem{listening} SC203 Group. \textit{SonicMaster Listening Study: Results from the First 100 Reviewed Samples}. 15 November 2026.
\end{thebibliography}

\end{document}

